% Source: Wald GR, physical PDF pp. 450--452 (printed pp. 437--439).
% Coverage: Appendix C.1, Maps of Manifolds; equations (C.1.1)--(C.1.5).
% Figures/tables/footnotes: none.
% Status: source-reviewed and compile-clean.
% Uncertainties: none.

\subsection{Maps of Manifolds}\label{sec:C.1}

Let \(M\) and \(N\) be manifolds (not necessarily of the same dimension) and let
\(\phi:M\to N\) be a \(C^\infty\) map. In a natural manner, \(\phi\) \lq\lq pulls
back\rq\rq\ a function \(f:N\to\mathbb R\) on \(N\) to the function
\(f\circ\phi:M\to\mathbb R\) obtained by composing \(f\) with \(\phi\). Similarly,
in a natural way, \(\phi\) \lq\lq carries along\rq\rq\ tangent vectors at \(p\in M\)
to tangent vectors at \(\phi(p)\in N\)---i.e., it defines a map
\(\phi^*:T_pM\to T_{\phi(p)}N\)---as follows: For \(v\in T_pM\) we define
\(\phi^*v\in T_{\phi(p)}N\) by
\begin{equation}
  (\phi^*v)(f)=v(f\circ\phi)
  \tag{C.1.1}\label{eq:C.1.1}
\end{equation}
for all smooth \(f:N\to\mathbb R\), where we have dropped the vector indices on
\(v\) and \(\phi^*v\) since that notation is inconvenient here. It is easy to check
that \(\phi^*v\) satisfies the properties required of a tangent vector at \(\phi(p)\)
and thus Eq.~\eqref{eq:C.1.1} properly defines the map \(\phi^*\). Note that
\(\phi^*\) is linear and may be viewed as the \lq\lq derivative of \(\phi\)\rq\rq\
at \(p\). [The matrix of components of \(\phi^*\) in the coordinate bases of a
coordinate system \(\{x^\nu\}\) at \(p\) and a coordinate system \(\{y^\mu\}\) at
\(\phi(p)\) equals the Jacobian matrix of the map \(\phi\) between the coordinates,
i.e., \((\phi^*)^\mu{}_\nu=\partial y^\mu/\partial x^\nu\).] By the implicit
function theorem, \(\phi:M\to N\) will be one-to-one in a neighborhood of \(p\) if
\(\phi^*:T_pM\to T_{\phi(p)}N\) is one-to-one.

Similarly, we can use \(\phi\) to \lq\lq pull back\rq\rq\ covectors at \(\phi(p)\) to
covectors at \(p\). We define the map \(\phi_*:T_{\phi(p)}^\vee N\to T_p^\vee M\) by
requiring that for all \(v^a\in T_pM\),
\begin{equation}
  (\phi_*\mu)_a v^a=\mu_a(\phi^*v)^a.
  \tag{C.1.2}\label{eq:C.1.2}
\end{equation}
We can extend the action of \(\phi_*\) to map tensors of type \((0,l)\) at
\(\phi(p)\) to tensors of type \((0,l)\) at \(p\) by
\begin{equation}
  (\phi_*T)_{a_1\ldots a_l}(v_1)^{a_1}\cdots(v_l)^{a_l}
  =T_{a_1\ldots a_l}(\phi^*v_1)^{a_1}\cdots(\phi^*v_l)^{a_l}.
  \tag{C.1.3}\label{eq:C.1.3}
\end{equation}
Similarly, we can extend the action of \(\phi^*\) to map tensors of type \((k,0)\)
at \(p\) to tensors of type \((k,0)\) at \(\phi(p)\) by
\begin{equation}
  (\phi^*T)^{b_1\ldots b_k}(\mu_1)_{b_1}\cdots(\mu_k)_{b_k}
  =T^{b_1\ldots b_k}(\phi_*\mu_1)_{b_1}\cdots(\phi_*\mu_k)_{b_k}.
  \tag{C.1.4}\label{eq:C.1.4}
\end{equation}
(By Eq.~\eqref{eq:C.1.2}, this is consistent with our original definition of
\(\phi^*\) on vectors.) However, in general we cannot extend \(\phi^*\) or
\(\phi_*\) to mixed tensors since \(\phi^*\) does not know how to \lq\lq carry
along\rq\rq\ lower-index tensors, while \(\phi_*\) does not know how to
\lq\lq pull back\rq\rq\ upper-index tensors.

As defined in Chapter~2, a \(C^\infty\) map \(\phi:M\to N\) is said to be a
\emph{diffeomorphism} if it is one-to-one, onto, and its inverse is \(C^\infty\). If
\(\phi\) is a diffeomorphism (which necessarily implies \(\dim M=\dim N\)), then we
can use \(\phi^{-1}\) to extend the definition of \(\phi^*\) to tensors of all types
by using the fact that \((\phi^{-1})^*\) goes from \(T_{\phi(p)}N\) to \(T_pM\). If
\(T^{b_1\ldots b_k}{}_{a_1\ldots a_l}\) is a tensor of type \((k,l)\) at \(p\), we
define the tensor
\((\phi^*T)^{b_1\ldots b_k}{}_{a_1\ldots a_l}\) at \(\phi(p)\) by
\begin{equation}
\begin{aligned}
  &(\phi^*T)^{b_1\ldots b_k}{}_{a_1\ldots a_l}
  (\mu_1)_{b_1}\cdots(\mu_k)_{b_k}
  (t_1)^{a_1}\cdots(t_l)^{a_l} \\
  &\qquad =
  T^{b_1\ldots b_k}{}_{a_1\ldots a_l}
  (\phi_*\mu_1)_{b_1}\cdots(\phi_*\mu_k)_{b_k}
  ([\phi^{-1}]^*t_1)^{a_1}\cdots([\phi^{-1}]^*t_l)^{a_l}.
\end{aligned}
  \tag{C.1.5}\label{eq:C.1.5}
\end{equation}
Similarly, we could extend the map \(\phi_*\) to all tensors. However, it is not
difficult to show that \(\phi_*=(\phi^{-1})^*\), so we need only consider
\(\phi^*\) and \((\phi^{-1})^*\).

If \(\phi:M\to M\) is a diffeomorphism and \(T\) is a tensor field on \(M\), we can
compare \(T\) with \(\phi^*T\). If \(\phi^*T=T\), then even though we have
\lq\lq moved \(T\)\rq\rq\ via \(\phi\), it has \lq\lq stayed the same.\rq\rq\ In
other words, \(\phi\) is a \emph{symmetry transformation} for the tensor field \(T\).
In the case of the metric \(g_{ab}\), a symmetry transformation---i.e., a
diffeomorphism \(\phi\) such that \((\phi^*g)_{ab}=g_{ab}\)---is called an
\emph{isometry}.

We have already remarked in Chapter~2 that if \(\phi:M\to N\) is a diffeomorphism,
then \(M\) and \(N\) have identical manifold structure. If a theory describes nature
in terms of a spacetime manifold, \(M\), and tensor fields, \(T^{(i)}\), defined on
the manifold, then if \(\phi:M\to N\) is a diffeomorphism, the solutions
\((M,T^{(i)})\) and \((N,\phi^*T^{(i)})\) have physically identical properties. Any
physically meaningful statement about \((M,T^{(i)})\) will hold with equal validity
for \((N,\phi^*T^{(i)})\). On the other hand, if \((N,T'^{(i)})\) is not related to
\((M,T^{(i)})\) by a diffeomorphism and if the tensor fields \(T^{(i)}\) represent
measurable quantities, then \((N,T'^{(i)})\) will be physically distinguishable from
\((M,T^{(i)})\). Thus, the diffeomorphisms comprise the gauge freedom of any theory
formulated in terms of tensor fields on a spacetime manifold. In particular,
diffeomorphisms comprise the gauge freedom of general relativity.

It is worth noting that an alternative viewpoint on diffeomorphisms can be taken.
Above, we have discussed diffeomorphisms without introducing or making any
reference to coordinate systems. We have taken an \lq\lq active\rq\rq\ point of view
by associating with \(\phi\) a map from tensors at \(p\) to tensors at \(\phi(p)\).
However, if we are given a coordinate system \(\{x^\mu\}\) covering a neighborhood,
\(U\), of \(p\) and a coordinate system \(\{y^\mu\}\) covering a neighborhood,
\(V\), of \(\phi(p)\), we may take the following \lq\lq passive\rq\rq\ point of view.
We may use \(\phi\) to define a new coordinate system \(x'^\mu\) in the neighborhood
\(O=\phi^{-1}[V]\) of \(p\) by setting \(x'^\mu(q)=y^\mu(\phi(q))\) for \(q\in O\).
We then may view the effect of \(\phi\) as leaving \(p\) and all tensors at \(p\)
unchanged, but inducing the coordinate transformation \(x^\mu\to x'^\mu\). This
\lq\lq passive\rq\rq\ point of view on diffeomorphisms is, philosophically,
drastically different from the above \lq\lq active\rq\rq\ viewpoint, but, in
practice, these viewpoints are really equivalent since the components of the tensor
\(\phi^*T\) at \(\phi(p)\) in the coordinate system \(\{y^\mu\}\) in the active
viewpoint are precisely the components of \(T\) at \(p\) in the coordinate system
\(\{x'^\mu\}\) in the passive viewpoint.
